\chapter*{学生实验}\markboth{学生实验}{学生实验}
\stepcounter{chapter}
\ctexset{section/name={实验,}}
\addcontentsline{toc}{chapter}{学生实验}
\renewcommand{\thefigure}{\arabic{figure}}
\renewcommand{\thetable}{\arabic{table}}

\section{化学实验基本操作}
\subsection[托盘天平的使用]{托盘天平的使用\texorpdfstring{\protect\footnote{教师可以根据实际情况选做。}}{}}
托盘天平（\cref{fig:a1}）是常用的称重器具，用于精确度不高的称量，一般能称准到 \qty{0.1}{g}（有的是 \qty{0.2}{g}）。
在初中化学课本的 “化学实验基本操作” 中已介绍过它的使用方法。

\begin{figure}
  \includegraphics{A-1.pdf}
  \caption{托盘天平}\label{fig:a1}
\end{figure}

\subsubsection*{学生练习}
\begin{enumerate}[1.]
  \item 称一个小烧杯的质量。烧杯应放在托盘天平的左盘还是右盘？为什么？
  \item 称取 \qty{2}{g} 食盐。在托盘天平的两盘各放一张同样的纸，为什么？注意当食盐的称取量只缺很少时，左手拿药匙，用右手轻拍左手手腕，小心振动药匙加足药量。
\end{enumerate}

\subsection{容量瓶的使用}
容量瓶是配制准确浓度溶液的仪器。
容量瓶是细颈、梨形的平底玻璃瓶（\cref{fig:a2}）。
瓶口配有磨口玻璃塞或塑料塞，它的颈部刻有标线，瓶上标有温度和容量。
常用的容量瓶有 \qtylist{50;100;250}{ml} 等几种。
\subsubsection*{容量瓶的使用方法}
\begin{enumerate}[1.]
  \item 使用容量瓶前检查它是否漏水\quad 方法如下: 往瓶内加水，塞好瓶塞，一手拿瓶，一手顶住瓶塞，把瓶倒立过来，观察瓶塞周围是否有水漏出。如果不漏水，把瓶塞旋转 \ang{180} 后塞紧，仍把瓶倒立过来，再检验是否漏水，经检查不漏水的容量瓶才能使用。
  \begin{figure}
    \begin{minipage}[b]{0.34\linewidth}\centering
      \includegraphics{A-2.pdf}
      \caption{容量瓶}\label{fig:a2}
    \end{minipage}
    \begin{minipage}[b]{0.34\linewidth}\centering
      \includegraphics{A-3.pdf}
      \caption{溶液从烧杯转移入容量瓶}\label{fig:a3}
    \end{minipage}
    \begin{minipage}[b]{0.3\linewidth}\centering
      \includegraphics{A-4.pdf}
      \caption{容量瓶的拿法}\label{fig:a4}
    \end{minipage}
  \end{figure}
  \item 配置溶液
  \begin{enumerate}[(1)]
    \item 如果试样是固体，把称好的试样溶解在烧杯里; 如果试样是液体，需用量筒量取倒入烧杯里，然后再加少量蒸馏水，用玻璃棒搅动，使它混和均匀。
    \item 把溶液从烧杯移到容量瓶里（\cref{fig:a3}）并多次洗涤烧杯，把洗涤液也移入容量瓶，以保证溶质全部移到容量瓶里。缓慢地加入蒸馏水，到接近标线 \qtyrange{2}{3}{cm} 处，用滴管滴加蒸馏水到标线（小心操作，切勿超过标线）。
    \item 盖好瓶塞，用食指顶住瓶塞，用另一只手的手指顶住瓶底（\cref{fig:a4}），把容量瓶倒转和摇动多次，使溶液混和均匀。
  \end{enumerate}
\end{enumerate}

\subsubsection*{学生练习}
练习容量瓶的使用方法（可以用水代替溶液）。

\subsection{萃取和分液操作}
\begin{enumerate}[1.]
  \item 萃取\quad 利用溶质在互不相溶的溶剂里溶解度的不同，用一种溶剂把溶质从它与另一溶剂所组成的溶液里提取出来的方法，叫做萃取。 苯取在生产和科学实验上用途很广。
  \item 分液\quad 分液是把两种不相混溶的液体分开的操作，使用的仪器是分液漏斗。分液漏斗有圆筒形、圆球形和圆锥形（\cref{fig:a5}）等几种，容积有 \qtylist{50;100;250}{ml} 等几种。萃取和分液有时可结合进行。
  \begin{figure}
    \begin{minipage}[b]{0.32\linewidth}\centering
      \includegraphics{A-5.pdf}
      \caption{分液漏斗}\label{fig:a5}
    \end{minipage}
    \begin{minipage}[b]{0.32\linewidth}\centering
      \includegraphics{A-6.pdf}
      \caption{倒转分液漏斗}\label{fig:a6}
    \end{minipage}
    \begin{minipage}[b]{0.32\linewidth}\centering
      \includegraphics{A-7.pdf}
      \caption{萃取操作}\label{fig:a7}
    \end{minipage}
  \end{figure}
  \item 操作方法（萃取和分液结合进行）
  \begin{enumerate}[(1)]
    \item 在溶液中加入萃取剂，用右手压住分液漏斗口部，左手握住活塞部分，把分液漏斗倒转过来用力振荡 (图 6)。
    \item 把分液漏斗放在铁架台上（\cref{fig:a7}），静置片刻。
    \item 把分液漏斗上的玻璃塞打开或使塞上的凹槽或小孔对准漏斗口上的小孔，使漏斗内外空气相通，以保证漏斗里的液体能够流出。
    \item 打开活塞，使下层液体慢慢流出。
  \end{enumerate}
\end{enumerate}

\subsubsection*{学生练习}
用量筒取 \qty{10}{ml} 碘的饱和水溶液，倒入分液漏斗，然后注入 \qty{3}{ml} 四氯化碳，振荡后，静置一会儿，溶液分成两层，上层是水层，下层是四氯化碳层。
打开活塞，用小烧杯盛接下层液体，倒入指定容器回收。


\section{配置一定摩尔浓度的溶液}
\subsection*{实验目的}
\begin{enumerate}[1.]
  \item 初步学会配制一定摩尔浓度溶液的方法；
  \item 初步学会容量瓶的使用和学习腐蚀性药品的称量。
\end{enumerate}

\subsection*{实验用品}
托盘天平和砝码、烧杯、量筒、玻璃棒、\qty{250}{ml} 容量瓶、胶头滴管、药匙。

浓盐酸（密度 \qty{1.19}{g/cm^3}、百分比浓度 37.5\%）、氢氧化钠。
\subsection*{实验步骤}
\begin{enumerate}[1.]
  \item 配制 \qty{250}{ml} $1M$ 盐酸
  \begin{enumerate}[(1)]
    \item 计算溶质的量\quad 根据浓盐酸密度（\qty{1.19}{g/cm^3}）、 百分比浓度 37.5\%，计算配制 \qty{250}{ml} $1M$ 盐酸需浓盐酸的毫升数。
    \item 用量筒量取浓盐酸\quad 用量筒量取所需的浓盐酸，倒入烧杯中，然后再加入少量水（约 \qty{50}{ml}），用玻璃棒慢慢搅动，使混和均匀，并冷却。
    \item 配制溶液\quad 把已稀释的盐酸沿玻璃棒注入容量瓶，并用 \qty{30}{ml} 水洗涤烧杯两次，洗液也注入容量瓶。振荡，使溶液混和均匀，然后继续往容量瓶中小心地加水，到液面接近刻度 \qtyrange{2}{3}{cm} 处时，改用胶头滴管加水到刻度线。把容量瓶盖紧，再振荡摇匀。这样得到的溶液就是 $1M$ 盐酸。
  \end{enumerate}
  \item 配制 \qty{250}{ml} $1.1M$ 的氢氧化钠溶液
  \begin{enumerate}[(1)]
    \item 计算溶质的量\quad 计算配制 \qty{250}{ml} $1.1M$ 的氢氧化钠溶液所需氢氧化钠的质量。
    \item 称量氢氧化钠\quad 用托盘天平先称量一干燥而洁净的烧杯的质量。然后把氢氧化钠放入烧杯，再称出它们的总质量。从总质量减去烧杯的质量便等于所需的氢氧化钠的质量。
    \item 配制溶液\quad 往烧杯中加入 \qty{50}{ml} 水，用玻璃棒搅动，使它溶解，并冷却，然后按照配制盐酸的方法配成 \qty{250}{ml} 的 $1.1M$ 的氢氧化钠溶液。
  \end{enumerate}
  把上面配成的溶液倒入指定的容器里。 
\end{enumerate}

\begin{Theorem}{讨论}
把烧杯里的溶液转移到容量瓶里后，进行洗涤，烧杯的洗涤液为什么也要倒入容量瓶？
\end{Theorem}

\section{重结晶法提纯硫酸铜\texorpdfstring{\quad}{ }测定硫酸铜晶体里结晶水的含量}
\subsection*{实验目的}
\begin{enumerate}[1.]
  \item 学习用重结晶法提纯晶体；
  \item 学习测定晶体里结晶水含量的方法；
  \item 初步掌握灼烧的技能。
\end{enumerate}
\subsection*{实验用品}
托盘天平和砝码、烧杯、漏斗、滤纸、量筒、玻璃棒、铁架台、石棉网、表面皿、瓷坩埚、坩埚钳、干燥器、酒精灯。

硫酸铜晶体。

\subsection*{实验步骤}
\begin{enumerate}[1.]
  \item 重结晶法提纯硫酸铜晶体
  
  为了得到纯度更高的晶体，可把结晶出来的晶体重新溶解在蒸馏水里，加热，制成饱和溶液，冷却，使它再一次结晶，杂质留在母液里。这就是重结晶或再结晶。
  \begin{enumerate}[(1)]
    \item 称量和溶解\quad 用托盘天平称取粗硫酸铜晶体 \qty{5}{g}，放在洁净的小烧杯里，然后用量筒量取 \qty{10}{ml} 蒸馏水倒入烧杯并加热，使硫酸铜完全溶解。
    \item 过滤\quad 趁热把溶液倾入事先组装好的过滤器的漏斗中，用一个小烧杯接滤液。
    \item 蒸发\quad 把烧杯放在石棉网上加热，蒸去三分之一体积的溶液。
    \item 结晶\quad 把烧杯浸到冷水里，溶液中即有硫酸铜晶体析出。
    \item 干燥\quad 小心倾倒出烧杯内的母液（回收），把晶体放在表面皿上、用滤纸吸收晶体表面的水，再把晶体放在两层滤纸上，用玻璃棒铺开，上面再盖一张滤纸，用手指轻轻压挤以吸去晶体表面的水。换新的滤纸，重复操作一次或两次，到晶体干燥为止。
  \end{enumerate}
  \item 硫酸铜晶体结晶水含量的测定
  \begin{enumerate}[(1)]
    \item 称量\quad 准确称量干燥的瓷坩埚的质量，并用这个坩埚称取 \qty{2}{g} 自己制备的硫酸铜晶体。
    \item 加热\quad 把盛有硫酸铜晶体的瓷坩埚放在石棉网上，用酒精灯慢慢地加热，直到硫酸铜晶体的蓝色完全变白，且不再逸出水蒸气为止。然后把坩埚放到干燥器里冷却。
    \item 称量\quad 待瓷坩埚在干燥器里冷却后，放在天平上称量，记下瓷坩埚和无水硫酸铜的质量。
    \item 加热再称量，到质量不再变\quad 把盛有无水硫酸铜的瓷坩埚再加热，放在干燥器里冷却后再称量，记下质量。到两次称量的误差不超过 \qty{0.1}{g} 为止。
    \item 计算\quad 根据实验结果，求 \qty{1}{mol} 硫酸铜晶体中含几摩尔结晶水，可通过下式进行计算：
    \[ \frac{W_{\ce{CuSO4}}}{M_{\ce{CuSO4}}}:\frac{W_{\ce{H2O}}}{M_{\ce{H2O}}} =1:x \]
    \begin{EqDesc}[公式中，]{W_{\ce{CuSO4}}}
      \item[x] 即是 \qty{1}{mol} 硫酸铜晶体中所含结晶水的摩尔数。
      \item[W_{\ce{CuSO4}}] 灼烧后无水硫酸铜的质量
      \item[M_{\ce{CuSO4}}] 硫酸铜的摩尔质量
      \item[W_{\ce{H2O}}] 晶体中所含结晶水的质量
      \item[M_{\ce{H2O}}] 水的摩尔质量
    \end{EqDesc}
  \end{enumerate}
\end{enumerate}
\begin{Theorem}{问题和讨论}
用重结晶法提纯硫酸铜晶体，在蒸发滤液时，为什么不能把滤液蒸干？
\end{Theorem}


\section{氯、溴、碘的性质}
\subsection*{实验目的}
\begin{enumerate}[1.]
  \item 认识氯、溴和碘以及卤化物的反应；
  \item 学习萃取和分液的操作方法；
  \item 认识卤素间的置换反应。
\end{enumerate}
\subsection*{实验用品}
试管、胶头滴管、分液漏斗、烧杯、表面皿、酒精灯、铁架台。

氯化钠、溴化钠、碘化钾、氯水、溴水、碘水、淀粉溶液、碘、四氯化碳、碘化钾淀粉试纸、硝酸银溶液、稀硝酸。
\subsection*{实验步骤}
\begin{enumerate}[1.]
  \item 观察氯水的颜色\quad 把贮氯水的瓶的瓶盖打开，小心地按正确操作的方法闻氯气的气味。
  \item 碘跟淀粉的反应\quad 在两个试管里，各加入淀粉溶液；然后向其中的一个试管里加 \numrange{2}{3} 滴碘水，向另一个试管加 \numrange{2}{3} 滴碘化钾溶液。观察到什么现象？并解释原因。
  \item 碘的升华\quad 在干燥的小烧杯里放几粒碘，盖上盛有少量冷水的表面皿，微微加热（\cref{fig:a8}）。观察碘晶体受热后所起的变化，解释发生的现象。回收凝结在表面皿底上的碘。
  \begin{figure}
    \includegraphics{A-8.pdf}
    \caption{碘的升华}\label{fig:a8}
  \end{figure}
  \item 萃取
  \begin{enumerate}[(1)]
    \item 在试管里加入一滴溴水，再加 5 滴四氯化碳，振荡试管，观察四氯化碳层颜色的变化。
    \item 用量筒取 \qty{10}{ml} 碘水，用碘化钾淀粉试纸试之。把碘水倒入分液漏斗，然后再加入 \qty{3}{ml} 四氯化碳，振荡，静置后分液，用小烧杯接四氯化碳溶液，回收。再用碘化钾淀粉试纸蘸取萃取后的碘水，和没萃取前的碘水进行比较。\footnote{如果碘化钾淀粉试纸仍然变蓝，可以再萃取一次。}
  \end{enumerate}
  \item 氯、溴、碘之间的置换反应
  \begin{enumerate}[(1)]
    \item 混和少量碘化钾溶液和淀粉溶液，分装在两个试管里，向一个试管里加入氯水，向另一个试管里加入溴水。各有什么现象发生：解释发生的现象。写出反应的化学方程式。
    \item 在两个试管里分别加入溴化钠溶液，然后向其中一个试管里加氯水，向另一个试管里加碘水。各有什么现象发生？ 解释发生的现象，并写出有关反应的化学方程式。
  \end{enumerate}
  \item 卤化物跟硝酸银的反应\quad 取三个试管，分别注入少量氯化钠溶液、溴化钠溶液、碘化钾溶液，然后在每个试管里各加入 \numrange{2}{3} 滴硝酸银溶液。再分别向每个试管里加入少量稀硝酸。观察发生的现象，加以解释，并写出有关反应的化学方程式。
\end{enumerate}
\begin{Theorem}{问题和讨论}
  碘化钾淀粉试纸常用来检验氯气，能不能用这种试纸来检验氯化钠或氯化钾里的氯，为什么？
\end{Theorem}


\section{硫酸的性质\texorpdfstring{\quad}{ }硫酸根离子的检验}
\subsection*{实验目的}
\begin{enumerate}[1.]
  \item 学习实验室制二氧化硫的方法；
  \item 认识硫酸的特性；
  \item 学会检验硫酸根离子的方法。
\end{enumerate}

\subsection*{实验用品}
烧瓶、单孔橡皮塞、导管、棉花、水槽、试管、玻璃棒、酒精灯、玻璃片。

亚硫酸钠、品红溶液、蓝色石蕊试纸、浓硫酸、盐酸、氯化钡溶液、硫酸钠溶液、碳酸钠溶液、铜片。

\subsection*{实验步骤}
\begin{enumerate}[1.]
  \item 二氧化硫的制取和性质
  \begin{enumerate}[(1)]
    \item 照\cref{fig:a9} 装配好仪器。检查它的气密性。
    \item 往烧瓶里放入少量的亚硫酸钠，再加入少量的浓硫酸，用带有导管的塞子立刻塞住烧瓶口，把导管的另一端插入一个干燥的试管里，管口用一小团棉花堵住。等一会儿把导管拔出插入盛有品红溶液的试管里。观察发生的现象，写出制取二氧化硫反应的化学方程式。
    \begin{figure}
      \begin{minipage}[b]{0.48\linewidth}\centering
        \includegraphics{A-9.pdf}
        \caption{制取二氧化硫的装置}\label{fig:a9}
      \end{minipage}
      \begin{minipage}[b]{0.48\linewidth}\centering
        \includegraphics{A-10.pdf}
        \caption{二氧化硫在水中的溶解}\label{fig:a10}
      \end{minipage}
    \end{figure}
    \item 把盛有二氧化硫的试管用拇指堵住，把试管倒立在盛水的水槽里（\cref{fig:a10}），移开拇指，观察发生的现象。等水不再上升后，在水里用拇指堵住管口拿出，正立过来，用蓝色石蕊试纸检验。解释发生的现象。
  \end{enumerate}
  \item 浓硫酸的特性
  \begin{enumerate}[(1)]
    \item\label{itm:expr_hso_1} 浓硫酸的稀释\quad 取一个试管，往其中注入约 \qty{5}{ml} 蒸馏水，然后小心地沿试管壁加入约 \qty{1}{ml} 浓硫酸。轻轻振荡后，用手触摸试管外壁，有什么现象发生？这溶液留作下面实验使用。
    \item 浓硫酸遇纸发生的现象\quad 用玻璃棒蘸取浓硫酸在纸（下面垫上玻璃片）上写字，观察字迹的变化，为什么？
    \item 在一试管中加入一小片铜片，并倒入实验 \ref{itm:expr_hso_1} 制得的稀硫酸 \qty{3}{ml}，观察有没有反应发生。在酒精灯上加热片刻，有没有反应发生，为什么？
    \item 再向另一试管中加入一小片铜片，并倒入 \qty{2}{ml} 浓硫酸，在酒精灯上小心加热（注意试管口不要对着任何人），并用润湿的蓝色石蕊试纸在试管口（注意不要触及试管口）检验所生成的气体。观察有什么现象发生？片刻后，停止加热，待试管中液体冷却后，把这些溶液沿试管壁倒入另一盛有 \qty{5}{ml} 水的试管中，观察溶液的颜色。
  \end{enumerate}
  \item 硫酸根离子的检验
  \begin{enumerate}[(1)]
    \item 在盛有稀硫酸溶液的试管里，滴加少量的氯化钡溶液，有什么现象发生？再向该试管中加入少量盐酸，有没有变化？
    \item 在两个试管中，分别加入少量硫酸钠溶液和碳酸钠溶液，并分别滴入少量氯化钡溶液，注意观察发生的现象。再向试管中分别加入少量盐酸，进一步观察发生的现象，解释这些现象。
  \end{enumerate}

  写出它们反应的化学方程式。
\end{enumerate}

\section{实验习题}
\subsection*{实验目的}
\begin{enumerate}[1.]
  \item 巩固已学各章的知识内容和实验技能；
  \item 复习非金属离子和酸根离子的鉴别方法。
\end{enumerate}
\subsection*{实验习题}
\begin{enumerate}[1.]
  \item 怎样用化学方法除去热水瓶内的水垢？
  \item 用实验证明盐酸里含有氢和氯两种元素？
  \item 有三组无色溶液，怎样通过实验把它们鉴别出来？
  \begin{enumerate}[(1)]
    \item 硫酸铵溶液和氯化铵溶液，
    \item 硫酸钠溶液和亚硫酸钠溶液，
    \item 稀硫酸和稀盐酸。
  \end{enumerate}
  \item 在一个试管里倒入三分之一容积的溴水，然后加入一些锌粉，用玻璃棒搅动并微微加热（如溶液还没褪色，可再加入一些锌粉），到溶液变成无色。把这溶液分别倒在两个试管里: 往其中的一个试管加入氯水，往另一个试管里加入硝酸银溶液。观察发生的现象，并加以解释。\footnote{可根据情况选做。}
  \item 不另用其它试剂，鉴别以下两组溶液。
  \begin{enumerate}[(1)]
    \item \ce{HCl}、\ce{CaCl2}、\ce{Na2CO3}。
    \item \ce{NaBr}、\ce{AgNO3}、\ce{HCl}、\ce{Na2CO3}。\footnote{可根据情况选做。}
  \end{enumerate}
\end{enumerate}

\section{碱金属及其化合物的性质}
\subsection*{实验目的}
\begin{enumerate}[1.]
  \item 巩固对碱金属及其化合物性质的认识；
  \item 学习从草木灰中提取钾盐的方法；
  \item 学习用焰色反应检验碱金属离子。
\end{enumerate}
\subsection*{实验用品}
烧杯、试管、镊子、玻璃片、小刀、铝箔、药匙、木条、漏斗、滤纸、酒精灯、玻璃棒、导管、橡皮塞、铁架台（带铁夹）、蓝色钴玻璃片、铂丝（或光洁无锈的铁丝）。

金属钠、碳酸钠、碳酸氢钠、碳酸钾、过氧化钠、石灰水、酚酞试液、草木灰、氯化钡溶液、盐酸、硝酸银溶液、稀硝酸。

\subsection*{实验步骤}
\begin{enumerate}[1.]
  \item 金属钠的性质
  \begin{enumerate}[(1)]
    \item 用镊子取出一小块金属钠，用滤纸把煤油擦干。把钠放在玻璃片上，用小刀切下绿豆大小的一块钠。注意钠的硬度，观察新切开钠的表面光泽。
    \item 在小烧杯里预先倒入一些水，然后把切下的钠用镊子放入烧杯里，并迅速用玻璃片将烧杯盖好。观察发生的现象，说明钠比水重还是比水轻。
    \item 另切一小块钠（绿豆大），用铝箔（事前用针刺一些小孔）包好，再用镊子夹住，放到\cref{fig:a11} 所示装置的试管口下。等试管中气体已收集满时，把试管倒着移近酒精灯点燃，有什么现象发生？ 说明反应中生成了什么气体。
    \begin{figure}
      \includegraphics{a-11.pdf}
      \caption{钠跟水起反应}\label{fig:a11}
    \end{figure}

    向烧杯里滴几滴酚酞试液，有什么现象发生？

    金属钠跟水发生了什么化学反应？ 写出反应的化学方程式。
  \end{enumerate}
  \item 过氧化钠的性质
  
  在一个干燥的试管里放入一小匙过氧化钠，观察过氧化钠的颜色。向试管里加入约 \qty{3}{ml} 水，并用带有火星的木条试验产生的气体，说明生成了什么物质。写出反应的化学方程式。
  \item 碳酸氢钠的性质
  
  在一干燥的试管里放入碳酸氢钠粉末，约占试管体积的 $1/6$。试管口用带有导管的塞子塞紧，并把试管用夹子固定在铁架台上，使管口略下倾（\cref{fig:a12}）。导管的一端浸在石灰水里。
  \begin{figure}
    \includegraphics{A-12.pdf}
    \caption{碳酸氢钠的受热分解}\label{fig:a12}
  \end{figure}

  加热碳酸氢钠，观察发生的现象。发生了什么反应，写出反应的化学方程式。当气泡已经很少时先把试管提高，使导管口露出石灰水面，移去装有石灰水的烧杯，再熄灭酒精灯。为什么要这样操作？
  \item 从草木灰中提取钾盐
  \begin{enumerate}[(1)]
    \item 溶解\quad 在烧杯里放入适量草木灰\footnote{如果用 \qty{250}{ml} 的烧杯，放半烧杯（\qty{10}{g} 左右）草木灰。}，加水到高于灰表面 \qtyrange{1}{2}{cm} 处，用玻璃棒搅动（可以适当加热），加速草木灰中钾盐的溶解。
    \item 过滤\quad 准备一个过滤器，把烧杯中的草木灰连同浸液一起过滤。如果滤液浑浊，再过滤一次，直到滤液澄清。
    \item 蒸发\quad 把滤液倒入蒸发皿，然后把蒸发皿如\cref{fig:a13} 放置，加热。用玻璃棒不断搅动液体，防止液滴飞溅。当蒸发到只剩少量液体时，停止加热。
    \begin{figure}
      \includegraphics{A-13.pdf}
      \caption{蒸发}\label{fig:a13}
    \end{figure}
    \item 冷却\quad 静置片刻，可以见到有钾盐晶体出现。
    \item\label{itm:experiment-5} 碳酸根离子、硫酸根离子、氯离子的检验
    
    取制得的少量晶体，放入试管，加人蒸馏水溶解并把溶液分成三份，分装在三个试管里：
    \begin{itemize}
      \item 在第一个试管里加入盐酸，观察发生的现象。
      \item 在第二个试管里加入氯化钡溶液，再加入盐酸，观察发生的现象。
      \item 在第三个试管里加入硝酸银溶液，再加入稀硝酸，观察发生的现象。
    \end{itemize}
  \end{enumerate}
  写出 \ref{itm:experiment-5} 中有关反应的化学方程式。

  根据上述反应，可以验证溶液里含有哪些离子？
  \item 焰色反应
  \begin{enumerate}[(1)]
    \item 准备碳酸钠溶液、碳酸钾溶液（或用它们的粉末）和碳酸钠、碳酸钾的混和溶液（或混和粉末）各少许。把铂丝烧热\footnote{如果没有铂丝，可以用光洁无锈的铁丝代替。}，然后用铂丝蘸一些碳酸钾的溶液（或粘一些碳酸钾的粉末），再放到酒精灯上灼烧，隔着蓝色钴玻璃观察火焰的颜色。火焰显出了什么颜色？为什么？
    \item 把铂丝洗干净，先后分别蘸碳酸钠溶液（或粉末），碳酸钾、碳酸钠混和溶液（或混和粉末），放到酒精灯上灼烧。在后一个实验里，先直接观察火焰颜色，再隔着蓝色钴玻璃观察，火焰各呈什么颜色？为什么？
    \item 把铂丝洗干净，蘸自己制得的钾盐颗粒，放到酒精灯上灼烧，隔着蓝色钴玻璃观察火焰的颜色。
  \end{enumerate}
\end{enumerate}
\begin{Theorem}{问题和讨论}
如果氯化钠里混有碳酸钠，怎样检验并除去？
\end{Theorem}


\section{同周期、同主族元素性质的递变}
\subsection*{实验目的}
巩固对同周期、同主族元素性质递变规律的认识。

\subsection*{实验用品}
试管、试管夹、酒精灯、烧杯、砂纸、药匙、玻璃片。

钾、钠、镁条、铝片、氢氧化钠溶液、酚酞试液、氯化镁溶液、氯化铝溶液、氢硫酸、氯水、溴水、氯化钠、溴化钠、碘化钠。

\subsection*{实验步骤}
\begin{enumerate}[1.]
  \item\label{itm:exper_step_1} 同周期元素性质的递变
  \begin{enumerate}[(1)]
    \item\label{itm:exper_step_1_1} 取一个 \qty{100}{ml} 的烧杯，向烧杯中注入约 \qty{50}{ml} 水，另取两个试管各注入约 \qty{5}{ml} 水，然后取绿豆大的一块钠，放入烧杯中。取一条镁条，用砂纸擦去表面的氧化物后，放入一个试管中。再取一片铝片，浸入氢氧化钠溶液中，以除去表面的氧化膜，然后取出，用水洗净，放入另一个试管中。注意观察反应现象，如果反应缓慢，可在酒精灯上加热。
    \item 向上述烧杯和两个试管中，各滴入 \numrange{2}{3} 滴酚酞试液，观察发生的现象。
    \item 取两个试管，分别加入 \qty{3}{ml} 氯化镁溶液和 \qty{3}{ml} 氯化铝溶液。然后逐渐滴入过量的氢氧化钠溶液，观察发生的现象。
    \item 在试管中加入约 \qty{3}{ml} 氢硫酸，然后滴入氯水，观察发生的现象。
  \end{enumerate}
  \par\medskip
  根据上述实验可得出什么结论？
  \item 同主族元素性质的递变
  \begin{enumerate}[(1)]
    \item 在一个 \qty{100}{ml} 的烧杯中，加人约 \qty{50}{ml} 水，然后加入一块绿豆大的钾，并用玻璃片把烧杯盖住。注意观察反应的剧烈程度，并与实验步骤 \ref{itm:exper_step_1}\ref{itm:exper_step_1_1} 中钠跟水的反应作比较。
    \item 在三个试管中，分别加入少量氯化钠、溴化钠、碘化钠晶体，并且各加入少量蒸馏水，使其溶解。然后分别加入 \qty{1}{ml} 氯水，注意观察溶液颜色的变化。
    \item 另取三个试管，用溴水代替上述实验中的氯水，做相同的实验。
  \end{enumerate}
  \par\medskip
  根据上述三个实验，可以得出什么结论？
\end{enumerate}
\begin{Theorem}{问题和讨论}
同一周期，从左到右元素的金属性逐渐减弱，非金属性逐渐增强; 同一主族，从上到下元素的金属性逐渐增强，非金属性逐渐减弱，这跟元素的原子结构和核外电子排布有什么关系？
\end{Theorem}

\section{实验习题}
\subsection*{实验目的}
\begin{enumerate}[1.]
  \item 巩固本书学过的知识内容和实验技能；
  \item 复习鉴别物质的方法。
\end{enumerate}
\subsection*{实验习题}
\begin{enumerate}[1.]
  \item 用实验方法鉴别碳酸钠和碳酸氢钠。
  \item 硫酸钠里混有少量硫酸铁的杂质，怎样把这杂质除去。
  \item 用实验方法实现下列变化：
  \[\ce{Ca(OH)2}\to\ce{CaCO3}\to\ce{Ca(HCO3)2}\to\ce{CaCO3}\to\ce{CaCl2}\]
  \item 用实验证明氯酸钾里含有钾、氧和氯三种元素。
  \item 有 \ce{NaCl}、\ce{NaBr}、\ce{NaI} 和稀盐酸四瓶无色溶液，通过实验把它们鉴别出来？
  \item 不另用其它试剂，鉴别以下溶液。\par \ce{H2SO4}、\ce{BaCl}、\ce{Na2SO4}、\ce{Na2CO3}。
  \item 只用一种试剂，鉴别以下物质。\footnote{此题可根据情况选做。}\par \ce{HCl}、\ce{NH4Cl}、\ce{Na2SO4}、\ce{(NH4)2SO4}。
\end{enumerate}

\section*{选作实验\texorpdfstring{\quad}{ }阿伏伽德罗常数的测定}
\addcontentsline{toc}{section}{选作实验\texorpdfstring{\quad}{ }阿伏伽德罗常数的测定}
\subsection*{实验目的}
\begin{enumerate}[1.]
  \item 明确认识阿佛加德罗常数的意义；
  \item 学习用单分子膜法测定阿佛加德罗常数的原理和操作。
\end{enumerate}
\subsection*{实验用品}
分析天平（\qty{0.001}{g}）、烧杯（\qty{50}{ml}）、玻璃棒、容量瓶 （\qty{250}{ml}）、胶头滴管、量筒（\qty{10}{ml}）、圆形水槽（直径 \qty{30}{cm}）、直尺。

硬脂酸（\ce{C17H35COOH}，分析纯）、苯（\ce{C6H6}，分析纯）、无水氯化钙。

\subsection*{实验步骤}
\begin{enumerate}[1.]
  \item 硬脂酸的苯溶液的制备\quad 教师于课前用小烧杯准确称取 \qtyrange{90}{100}{mg}硬脂酸\footnote{如果是托盘天平，称取硬脂酸 \qty{1.0}{g}，用苯溶解后，倒入 \qty{100}{ml} 容量瓶，加苯到刻度线。用吸量管吸取 \qty{3}{ml} 上面配制的溶液，放入 \qty{100}{ml} 容量瓶，再加苯到刻度线。}，以少量苯（\qty{300}{ml} 苯中加无水氯化钙  \qty{5}{g} ，充分振荡后过滤，以除去苯中含有的少量水分）溶解后，注入 \qty{250}{ml} 容量瓶中。用苯冲洗烧杯数次，洗液都倒入容量瓶。再往容量瓶中加苯到刻度线，摇匀即成。此溶液供学生实验使用。
  \item 测定从胶头滴管滴出的每滴硬脂酸的苯溶液的体积取一尖嘴拉得较细的胶头滴管，吸入硬脂酸的苯溶液，往小量筒中滴入 \qty{1}{ml} 然后记下它的滴数，并计算出一滴硬脂酸的苯溶液的体积 $v$。
  \item 测定水槽中水的表面积\quad 用直尺从三个不同方位准确量出水槽的内径，取其平均值。
  \item 硬脂酸单分子膜的形成\quad 用胶头滴管（如滴管外有溶液，用滤纸擦去）吸取硬脂酸溶液，在距水面约 \qty{5}{cm} 处，垂直往水面上滴一滴，待苯全部挥发，硬脂酸全部扩散至看不到油珠时，再滴第二滴。如此逐滴滴下，直到滴下一滴后，硬脂酸溶液不再扩散，而呈透镜状时为止。记下所滴硬脂酸溶液的滴数 $d$。
  \begin{figure}
    \includegraphics{a-14.pdf}
    \caption{单分子膜的形成}\label{fig:a14}
  \end{figure}
  \item 把水槽中水倒掉，用清水将水槽洗刷干净后，注入半槽水，重复以上操作二次。重复操作时，先挤净滴管内剩余的溶液，吸取新鲜溶液，以免由于滴管口的苯挥发引起溶液浓度的变化。取三次结果的平均值。
  \item 计算
  \begin{enumerate}[(1)]
    \item 如称取硬脂酸的质量为 $m$（\unit{g}），配成硬脂酸的苯溶液的体积为 $V$\,\unit{ml}，那么每毫升硬脂酸的苯溶液中含硬脂酸的质量为 $\dfrac{m}{V}$\,\unit{g}。
    \item \label{itm:experiment-calc2} 测得每滴硬脂酸苯溶液的体积为 $v$（\unit{ml}，形成单分子膜滴入硬脂酸溶液的滴数为 $d-1$\footnote{当最后一滴硬脂酸溶液滴下后，这滴溶液在水面呈透镜状，说明这滴溶液没有扩散，即没有参与单分子膜的形成。这时单分子膜已经形成完毕，应停止滴入溶液，所以，在计算形成单分子膜所需硬脂酸溶液的滴数时，应将最后一滴减掉，即滴数计为 $(d-1)$。}，那么形成单分子膜需用硬脂酸的质量（\unit{g}）为：
    \[v(d-1)\cdot\frac{m}{V}=\frac{vm(d-1)}{V} \]
    \item \label{itm:experiment-calc3} 根据水槽直径，计算出水槽中水的表面积 $S$（\unit{cm^2}）。实验测得每个硬脂酸分子的截面积 $A=\qty{2.2e-15}{cm^2}$，在水面形成的硬脂酸单分子膜中含有硬脂酸的分子个数为：$\dfrac{S}{A}$
    \item 根据 \ref{itm:experiment-calc2} 和 \ref{itm:experiment-calc3} 的结果，可计算出每个硬脂酸分子的质量（\unit{g}）为：
    \[\frac{\dfrac{vm(d-1)}{V}}{\dfrac{S}{A}}=\frac{Amv(d-1)}{SV}\]
    \item 每摩尔硬脂酸的质量等于 \qty{284}{g}，所以每摩尔硬脂酸中含有硬脂酸的分子个数，即阿佛加德罗常数 $N$ 为：
    \[ N=\frac{\mathrm{M}}{\dfrac{Amv(d-1)}{SV}}=\frac{\mathrm{M}SV}{Amv(d-1)}\]
  \end{enumerate}
\end{enumerate}
\begin{Theorem}{问题和讨论}
  采取哪些措施，可以减小这个实验的误差？
\end{Theorem}